\documentclass[aspectratio=169]{beamer}
\usetheme{Madrid}
\usecolortheme{beaver}
\usepackage[utf8]{inputenc}
\usepackage[T1]{fontenc}
\usepackage{tikz}
\usepackage{amssymb}
\usepackage{graphicx}
\setbeamertemplate{navigation symbols}{}
\setbeamercolor{title}{fg=white,bg=cyan!60!black}
\setbeamercolor{frametitle}{fg=white,bg=cyan!50!black}
\setbeamercolor{block title}{bg=cyan!15,fg=black}
\setbeamercolor{block body}{bg=cyan!5,fg=black}

\title{Story Studio 1 — Page Basics}
\subtitle{Lesson 4: Read More}
\author{ByteForge --- Grade 7 --- Game Maker}
\date{}

\begin{document}

\begin{frame}
    \titlepage
\end{frame}

\begin{frame}{A link connects one page to another}
    \Large
    Every link needs two things: the text someone clicks, and the address it goes to. Links are what turn separate pages into a connected site.
\end{frame}

\begin{frame}{A link connects one page to another --- Worked Example}
    \Large
    Worked example: a link block missing its destination is like a signpost with an arrow but no name on it — technically still a signpost, but useless to anyone reading it.
\end{frame}

\begin{frame}{Your Turn}
    \Large
    Open the student map, find \textbf{Story Studio 1 — Page Basics}, and complete
    the \textbf{Read More} challenge --- everything you just saw, they now do themselves.
\end{frame}

\end{document}
